% !TEX root = ../main.tex
% Caracterizacion del corpus (IT-58).
%
% Las cifras las imprime `scripts/verificadores/check_chunks.py` sobre
% `data/chunks.json`. Se copiaron de su salida del 29/08/2026, NO se estimaron.
% Si se regenera el corpus hay que volver a ejecutarlo y rehacer esta tabla: no
% ajustar un numero a mano, porque una diferencia puede ser un cambio legitimo
% de la fuente o una perdida silenciosa.
\begin{table}[htbp]
  \centering
  \caption{\textcolor{borrador2}{Caracterización de la colección documental
      sobre la que se han hecho los experimentos de la Fase~2.}}
  \label{tab:caracterizacion_corpus}
  \small
  \begin{tabular}{@{}l r@{}}
    \toprule
    \textbf{Magnitud}                              & \textbf{Valor}           \\
    \midrule
    Titulaciones                                   & 12                       \\
    \quad grados simples                           & 7                        \\
    \quad dobles grados                            & 5                        \\
    \midrule
    Asignaturas                                    & 528                      \\
    \quad con guía docente                         & 442                      \\
    \quad sin contenido de guía                    & 86                       \\
    \quad no ofertadas este curso                  & 10                       \\
    \quad sin créditos publicados                  & 1                        \\
    \midrule
    Guías docentes                                 & 288                      \\
    Bloques de salidas profesionales               & 8                        \\
    \midrule
    Fragmentos                                     & 1\,499                   \\
    Unidades documentales                          & 398                      \\
    \quad compartidas entre titulaciones           & 81                       \\
    \midrule
    Tamaño mínimo de fragmento (caracteres)        & 96                       \\
    Tamaño mediano                                 & 836                      \\
    Percentil 90                                   & 895                      \\
    Tamaño máximo                                  & 900                      \\
    Fragmentos por debajo del mínimo de 200        & 12 (0,80\,\%)            \\
    \bottomrule
  \end{tabular}
  \begin{minipage}{0.95\textwidth}
    \vspace{4pt}
    \tiny
    \textcolor{borrador2}{\textbf{Notas:} \textbf{Unidad documental}: la guía de
      una asignatura, el bloque de salidas de una titulación o la ficha de una
      asignatura sin guía; una misma unidad puede producir varios fragmentos y
      pertenecer a varias titulaciones. \textbf{Sin contenido de guía}: 81
      asignaturas cuya guía la Escuela no ha publicado y 5 que la publican con
      la plantilla vacía. \textbf{El mínimo de 200 caracteres es una
      preferencia, no una restricción}: de los 12 fragmentos que no lo alcanzan,
      2 son colas que no cabían junto a su vecino sin pasarse del máximo y 10
      son unidades enteras más cortas que el mínimo. Procedencia: extracción del
      16/08/2026, troceado del 19/08/2026, curso 2026-27; las tres viajan dentro
      de los propios ficheros. Cifras de \texttt{check\_chunks.py}.}
  \end{minipage}
\end{table}
